% SPDX-License-Identifier: GPL-3.0-or-later
% Copyright (C) 2026 Alessandro Greco
\chapter{流、SENDME 与拥塞控制}\label{cap:11-stream}
\aggiornatoal{2026 年 9 月 23 日}{torspec 提交 \texttt{8e1f245}；tor 0.4.9.13；提案 289、324}

\begin{obiettivi}
学完本章后，你将能够：
\begin{itemize}[nosep]
  \item 使用 BEGIN、CONNECTED、DATA 和 END 通过电路打开、使用和关闭 TCP 流；
  \item 解释为什么域名解析发生在出口，以及 RESOLVE 如何工作；
  \item 描述 Tor 原始的流量控制机制，包括电路窗口、流窗口和 SENDME；
  \item 理解为什么 SENDME 必须经过认证，以及它究竟证明了什么；
  \item 描述提案 324 引入的拥塞控制，以及 XON、XOFF 的作用；
  \item 通过自己构建的三跳电路下载文件，并测量真实出口的窗口行为（实验~5）。
\end{itemize}
\end{obiettivi}

\Cref{cap:10-cifratura-relay} 中的电路已经是一条通向出口的加密通道，但还没有承载真正有用的应用连接。本章让它开始工作：把客户端的 TCP 连接通过电路送往互联网，这些连接在 Tor 中称为\textbf{流}。随后我们处理每个网络都必须面对、而在 Tor 中尤其棘手的问题：如何防止发送方比网络或接收方能够消化的速度更快地发送数据，从而把中继队列塞满。

% ---------------------------------------------------------------------------
\section{流}\label{sec:stream}

流是客户端的一条 TCP 连接，例如到 Web 服务器的连接，通过某条电路承载。一条电路可以同时承载多个流；它们由 RELAY 消息中的 \textbf{StreamID} 区分（\cref{sec:relay}）。StreamID 由客户端选择，而且必须非零。\cref{fig:stream} 展示了完整生命周期。

\begin{figure}[htbp]
\centering
% SPDX-License-Identifier: GPL-3.0-or-later
% Copyright (C) 2026 Alessandro Greco
%% 简体中文版 EPUB 可访问性替代文本。
\altfig{客户端—出口—服务器数据流示意图：BEGIN 连接 example.com:80，可在 CONNECTED 之前发送乐观 DATA。出口单独执行 DNS 解析，然后通过 TCP SYN/SYN-ACK 建立连接，返回 CONNECTED。DATA 经加密 Tor 电路传到出口，再经外部 TCP 链路传到服务器。经典电路流控中的 SENDME 对应每 100 个携带 DATA 的信元增量，而不是 31 个。TCP FIN 后可能发送 END (DONE)。出口外的应用数据仍可由 HTTPS 保护。时序仅作示意。}
% 数据流生命周期
\begin{tikzpicture}[font=\sffamily\small,>=Stealth,
  msg/.style={font=\ttfamily\scriptsize,fill=white,inner sep=1pt},
  tcp/.style={-{Stealth[length=2.5mm]},thick,torgrigio}]
  \node[client] (c) at (0,0) {客户端};
  \node[relay,fill=torverde!30] (e) at (6.5,0) {出口};
  \node[server] (s) at (10.5,0) {Web 服务器};
  \foreach \x in {0,6.5,10.5} \draw[torgrigio] (\x,-0.5) -- (\x,-7.3);
  \draw[cellaf] (0,-1.0) -- node[msg,above]{BEGIN "example.com:80"} (6.5,-1.0);
  \draw[cellaf,dashed] (0,-1.6) -- node[msg,above]{DATA（乐观）} (6.5,-1.6);
  \draw[tcp] (6.5,-1.3) -- node[msg,above]{TCP SYN} (10.5,-1.3);
  \draw[tcp] (10.5,-2.1) -- node[msg,above]{SYN-ACK} (6.5,-2.1);
  \draw[risposta] (6.5,-2.6) -- node[msg,above]{CONNECTED（IP、TTL）} (0,-2.6);
  \draw[tcp] (6.5,-2.9) -- node[msg,above]{请求} (10.5,-2.9);
  \draw[tcp] (10.5,-3.6) -- node[msg,above]{响应} (6.5,-3.6);
  \draw[risposta] (6.5,-4.1) -- node[msg,above]{DATA, DATA, \ldots} (0,-4.1);
  \draw[cellaf] (0,-4.8) -- node[msg,above]{SENDME（电路，100 个 DATA 信元）} (6.5,-4.8);
  \draw[risposta] (6.5,-5.5) -- node[msg,above]{DATA, \ldots} (0,-5.5);
  \draw[tcp] (10.5,-6.1) -- node[msg,above]{FIN} (6.5,-6.1);
  \draw[risposta] (6.5,-6.7) -- node[msg,above]{END（DONE）} (0,-6.7);
  \node[font=\scriptsize,text=torgrigio] at (3.25,-7.6) {电路内部（三跳、已加密）};
  \node[font=\scriptsize,text=torgrigio] at (8.5,-7.6) {普通 TCP 连接};
\end{tikzpicture}

\caption{到 Web 服务器的一条流的生命周期。客户端与出口之间的消息通过电路并受到分层加密；出口与服务器之间则是普通 TCP 连接，如果应用本身没有使用 TLS，这一段不会自动加密。客户端可以在收到 CONNECTED 之前先发送数据，这称为\emph{乐观数据（optimistic data）}。}
\label{fig:stream}
\end{figure}

\subsection{打开流：BEGIN 与 CONNECTED}

客户端用一个新的 StreamID 向出口发送 \texttt{BEGIN}（RELAY 命令 1），以打开流。消息载荷指定目标~\spec{tor-spec}{Opening streams and transmitting data}：

\begin{center}
\begin{tabular}{lll}
\toprule
字段 & 长度 & 内容 \\
\midrule
ADDRPORT & 可变 & \texttt{address:port}，以零字节结尾 \\
FLAGS & 4（可选） & IPv4/IPv6 偏好 \\
\bottomrule
\end{tabular}
\end{center}

地址可以是 DNS 名称、IPv4 地址，或放在方括号中的 IPv6 地址。两个看似细小的规则体现了协议对不可区分性的重视：字符串必须以\textbf{小写}发送，避免大小写习惯暴露客户端实现；FLAGS 为零时必须\textbf{完全省略}。FLAGS 的第 1 位表示接受 IPv6，第 2 位表示拒绝 IPv4，第 3 位表示优先 IPv6。

收到 \texttt{BEGIN} 后，出口会：
\begin{enumerate}
  \item 检查自己不是这条电路的第一跳：电路到达出口的通道必须由共识中列出的、已认证中继打开。Tor 不能被当作单跳代理使用，否则出口的法律和滥用风险会更高；这样的尝试会用 \cella{DESTROY} 拒绝；
  \item 应用自身的\emph{出口策略（exit policy）}，也就是允许连接的地址与端口规则（\cref{cap:22-relay}）；
  \item 如有需要解析域名，然后建立 TCP 连接。
\end{enumerate}

成功后，出口回复 \texttt{CONNECTED}（RELAY 命令 4）。载荷通常含到达的 IPv4 地址（4 字节）以及客户端可以缓存它的秒数 TTL（4 字节）。IPv6 格式略有不同，也允许空载荷。客户端\emph{可以}忽略地址；如果缓存，则不得在其他电路复用，因为恶意出口可能故意返回一个唯一地址，以便后来重新识别该客户端。

\subsection{传输数据：DATA}

应用数据使用 \texttt{DATA}（RELAY 命令 2）传输；传统格式中每条最多承载 498 字节。出口把 DATA 的内容写入 TCP 连接，并把服务器返回的数据反向封装为 DATA。客户端不必等到 \texttt{CONNECTED} 才能开始发送：在 \texttt{BEGIN} 后可以立即发出 \texttt{DATA}，出口先把它们排队，连接成功后立即转发。这种\emph{乐观数据}省掉一次完整电路往返延迟，而 Tor 电路的 RTT 很容易达到数百毫秒。

\subsection{关闭流：END}

主动关闭连接或遇到错误的一方发送 \texttt{END}（RELAY 命令 3），载荷以 1 字节原因码开头~\spec{tor-spec}{Closing streams}。常见原因包括 \texttt{DONE}（6，正常结束）、\texttt{CONNECTREFUSED}（3）、\texttt{RESOLVEFAILED}（2）、\texttt{TIMEOUT}（7）和 \texttt{EXITPOLICY}（4）；最后一种后面还会带被拒绝的地址与 TTL。对于由自己发起的流，客户端应只发送 \texttt{MISC}（1）：与 \cella{DESTROY} 一样，没有必要向另一端暴露更多内部状态。

\subsection{在出口解析 DNS}\label{sec:stream-dns}

在 \texttt{BEGIN} 中，客户端可以直接提供\emph{域名}而不是 IP 地址，由出口负责解析。这不仅是便利性设计，更是安全要求。如果客户端在本地解析，它就会在 Tor 之外向自己的 DNS 解析器查询即将访问的每一个站点；任何观察这些查询的人——首先就是网络接入提供商——都能知道它准备访问什么。这就是常见的 \enquote{DNS 泄漏（DNS leak）}，练习~\ref{es:19-socks} 会演示它，实验~18（\cref{sec:lab-opsec}）则会从 tor 侧进行观察。

如果客户端只想查询地址而不建立连接，可以使用 \texttt{RESOLVE}（RELAY 命令 11），载荷是一个零终止的名称。出口以 \texttt{RESOLVED}（12）返回一个或多个“类型、长度、值、TTL”结构：\texttt{0x04} 表示 IPv4，\texttt{0x06} 表示 IPv6，\texttt{0x00} 表示名称（反向查询），\texttt{0xF0} 和 \texttt{0xF1} 分别表示临时和永久错误~\spec{tor-spec}{Remote hostname lookup}。tor 客户端能够通过本地 \opt{DNSPort} 向应用程序提供 DNS 服务，其底层正是这一机制。

\subsection{目录流：BEGIN\_DIR}

\file{BEGIN\_DIR}（RELAY 命令 13）不会把流打到互联网，而是连接中继自身的目录服务。客户端可以借此下载共识与描述符（\cref{cap:12-directory}），同时避免建立未加密的外部连接。它的载荷为空，而且出口策略不适用，因为流从未离开 tor 进程。

% ---------------------------------------------------------------------------
\section{原始流量控制}\label{sec:flusso-classico}

\subsection{问题}

设想出口的互联网连接很快，而客户端很慢。如果出口不断从服务器读取数据并全部塞进电路，中间中继的队列会持续增长，消耗内存并拖慢其他电路。TCP 能在每一条独立链路上解决拥塞，但一条 Tor 电路是多条不同 TCP 连接的串联：中间中继与客户端之间出现的 TCP 反向压力（backpressure），不会自动传播到出口。因此需要一种\emph{端到端}机制。

\subsection{窗口与 SENDME}

Tor 原始流量控制使用按 \texttt{DATA} 信元计数的\textbf{窗口}~\spec{tor-spec}{Flow control}：
\begin{itemize}
  \item \textbf{电路级}：发送方拥有初始为 1000 信元（或共识参数 \texttt{circwindow} 指定值）的\emph{发送窗口}。每发送一条 DATA 就减一，降到零必须停止。接收方维护对称的\emph{交付窗口}；每收到 100 条 DATA，就发送一条 StreamID 为零的 \texttt{SENDME}（RELAY 命令 5），让发送方窗口恢复 100。
  \item \textbf{流级}：窗口初始为 500 信元，每收到一条带有该流 StreamID 的 \texttt{SENDME}，恢复 50。
\end{itemize}

因为客户端可以直接与电路中的任意跳通信（\cref{sec:tor1}），所以会为每一跳维护一对窗口。这些窗口只统计由该端\emph{产生或消费}的 DATA，而不统计纯转发信元。

因此，在任意时刻每条电路最多有 1000 条 DATA 在途，每条流最多 500 条，约分别为 500\,KB 和 250\,KB。这样能够限制中继内存占用，却也限制吞吐量：若 RTT 为半秒，一条电路最多只能每半秒补满 1000 信元，也就是约 1\,MB/s，无论物理带宽多高。反过来，在拥塞网络中，1000 条在途信元又可能太多，队列和延迟都会增长。

\subsection{经过认证的 SENDME}

恶意客户端可以伪造 \texttt{SENDME}，即使根本没有真正接收数据，也不断告诉出口\enquote{可以继续发}。出口于是持续注入数据，这些数据在网络中积压，使客户端可以很低成本地耗尽中继内存或放大流量。从 0.4.1.1-alpha 开始，电路级 \texttt{SENDME} 因此需要\textbf{认证}（\proposta{289}）。载荷包含版本号（1）、长度（20）以及触发该 \texttt{SENDME} 的最后一条 \texttt{DATA} 信元对应的\textbf{累积摘要}。SENDME 的接收者在发送那条 DATA 时已经计算过该摘要，因此可以比较；不匹配就关闭电路。

摘要取决于此前收到的所有信元内容，但如果客户端预先知道数据，例如下载一份已知文件，它仍可能预测摘要。正因如此，每 100 个信元中至少必须有一条包含不可预测的随机字节；RELAY 消息中的随机填充（\cref{sec:relay}）也承担了这一作用。

\begin{approfondimento}[究竟是哪一个摘要？]
规范写的是\enquote{位于 SENDME 之前的信元摘要}。在 tor 实现中，它具体是反向方向在最后一条被识别信元之后的完整 20 字节 SHA-1 运行摘要，由 \file{src/core/crypto/relay\_crypto\_tor1.c} 中的 \file{tor1\_relay\_digest\_matches\_v0} 保存到 \file{sendme\_digest}。使用 CGO 时，则由最后一条信元的认证标签取代（\cref{sec:cgo}）。
\end{approfondimento}

% ---------------------------------------------------------------------------
\section{拥塞控制}\label{sec:cc}

固定窗口有两个相反问题：快速电路上太小，拥塞电路上又太大。Mike Perry 的 \proposta{324} 为 Tor 引入真正类似 TCP 的\textbf{拥塞控制}，从 tor 0.4.7.7（2022 年 4 月）开始部署。

\subsection{信号：往返时间}

出口怎样判断一条电路正在拥塞？答案是把 \texttt{SENDME} 当成时钟。出口记录发送那条将触发下一次 SENDME 的 \texttt{DATA} 的时间；SENDME 返回时，就得到电路的\textbf{往返时间}（RTT）。如果 RTT 相比观察过的最小值变大，说明沿途队列正在积压，发送方应减速。发送方据此调整\emph{拥塞窗口（cwnd）}，试图让它接近电路的带宽时延积，也就是在不产生额外排队时电路可以容纳的数据量。

\subsection{算法}

提案描述了三种受 TCP 启发的算法：Tor Westwood、Tor Vegas、Tor NOLA。经过 Shadow 仿真（\cref{cap:19-rete-privata}）以及真实网络评估后，tor 默认选择\textbf{Vegas}，对应 \file{src/core/or/congestion\_control\_common.c} 中的 \file{CC\_ALG\_DFLT}。与 TCP Vegas 类似，它通过当前窗口和最低观测 RTT 所需窗口之间的差异，估计有多少信元正处于队列中，再增减 cwnd，把估计排队量维持在目标阈值附近。

\subsection{参数：为什么是 31？}

拥塞控制通过 ntor-v3 对每一跳分别协商（\cref{sec:ntor3}）。客户端发送 \file{CC\_FIELD\_REQUEST}，中继回复 \file{CC\_FIELD\_RESPONSE}，其中包含 \file{sendme\_inc}：每收到多少条 \texttt{DATA} 后要求一条 \texttt{SENDME}。\cref{cap:09-circuiti} 中实验出口返回 31。这个值并非随意选择；tor 把它定义为一个最大 TLS 记录能容纳的信元数量：
\[
\file{TLS\_RECORD\_MAX\_CELLS}
= \left\lfloor \frac{16384}{514} \right\rfloor
= 31.
\]
初始窗口为 $4\times31=124$ 信元。与旧方案每 100 信元一次相比，每 31 信元一次 SENDME 能更频繁地测量 RTT，也能更快对拥塞变化作出反应。

\subsection{不再有流窗口：XON 与 XOFF}

启用拥塞控制后，流级 \texttt{SENDME} 消失。旧的流窗口存在两难：窗口小，会限制单个下载速度；窗口大，恶意客户端又可以同时打开许多流，让出口为每条流缓冲一个完整窗口。提案用两条消息取而代之：
\begin{description}
  \item[XOFF]（RELAY 命令 43）：\enquote{暂停在这条流上发送}；当面向应用程序的队列超过阈值时发出；
  \item[XON]（RELAY 命令 44）：\enquote{恢复发送}，并携带 4 字节字段说明该队列当前的排空速率，以便对端自适应。字段名为 \texttt{kbps\_ewma}，但 \proposta{324} 没有定义单位。tor 源码（\file{src/core/or/congestion\_control\_flow.c}）实际以\textbf{每秒 1000 字节}为单位计算和解释，而不是 1000 比特；变量名具有误导性，若按比特/秒解释就会产生八倍误差。
\end{description}

相比之下，电路级 \texttt{SENDME} 在信元被解密并识别后立即发送，发生在交给应用程序之前，因此测量的是\emph{网络}而不是应用处理速度。

\begin{attenzione}[拥塞控制与匿名性]
任何会根据流量作出反应的机制，都可能成为侧信道。观察窗口行为的攻击者可能推断电路内部状态。\proposta{324} 专门讨论这些风险和缓解措施，其中包括限制 XON/XOFF 的发送频率。这与 \cref{cap:08-celle} 中统一让 \cella{DESTROY} 使用原因码 NONE 属于同一种防御思路：任何依赖内部状态、又能被外部观察到的变化，都可能泄露信息。\cref{cap:26-attacchi} 会讨论利用流量线索的攻击。
\end{attenzione}

% ---------------------------------------------------------------------------
\section{实验：通过 Tor 下载文件}\label{sec:lab-stream}

\begin{laboratorio}[5]{流与流量控制}
\textbf{目标。} 编写 \file{labs/zh-Hans/代码/torlab/circuit.py}，把前面实验组合进 \texttt{Circuit} 类；再编写 \file{labs/zh-Hans/代码/torlab/stream.py} 打开流并传输数据。通过三跳电路从本地 Web 服务器下载文件，并测量出口窗口。\\
\textbf{前置条件。} 实验~1--4；正在运行的实验网络；本地 Web 服务器（\cmd{python3 -m http.server 8080}）。\\
\textbf{预计时间。} 3 小时。
\end{laboratorio}

\paragraph{步骤 1：一行构建电路。} \texttt{build} 打开到守卫的通道、建立第一跳，再逐步扩展；和 tor 一样，只在最终跳请求拥塞控制：

\begin{lstlisting}[language=Python]
circ = build([relay("003r"), relay("004r"), relay("006e")])
print(circ.exit_extensions)             # [(2, b'\x1f')]：sendme_inc = 31
\end{lstlisting}

\paragraph{步骤 2：打开流。}

\begin{lstlisting}[language=Python]
def begin_body(address, port, flag=0):
    body = f"{address.lower()}:{port}".encode() + b"\x00"
    return body + (struct.pack(">I", flag) if flag else b"")

s = Stream(circ, 1)
s.open("127.0.0.1", 8080)       # BEGIN，然后等待 CONNECTED 或 END
\end{lstlisting}

出口以 \texttt{CONNECTED} 回复 8 字节载荷：\texttt{7f000001 00000708}，即地址 127.0.0.1 与 TTL \texttt{0x708}=1800 秒。

\paragraph{步骤 3：SENDME。} 电路统计从出口收到的每一条 \texttt{DATA}，每累计 \file{sendme\_inc} 条就发送一条经过认证的 \texttt{SENDME}：

\begin{lstlisting}[language=Python]
def count_received_data(self):
    est = dict(self.exit_extensions)
    sendme_inc = est[2][0] if 2 in est else 100
    self.data_to_ack += 1
    if self.data_to_ack == sendme_inc:
        self.data_to_ack = 0
        digest = self.crypto.hops[-1].backward.digest.digest()
        body = struct.pack(">BH", 1, 20) + digest
        self.send(RelayMessage(RelayCommand.SENDME, 0, body))
\end{lstlisting}

\begin{attenzione}[一个很有教育意义的错误]
实验最初版本把这个计数器放在\emph{流}对象里。第一条流只返回两条信元，所以运行正常；第二条流下载 1\,MB 文件时，却被出口用 \cella{DESTROY} 中断。原因是电路级 \texttt{SENDME} 统计的是这条电路上\emph{所有流}的 \texttt{DATA}。出口已经把第一条流的两条信元算进去，因此期待第二条流的第 29 条信元后收到 SENDME；我们却等到第 31 条才发，而且摘要已不是出口期待的那个值，于是出口严格按照 \proposta{289} 关闭电路。把计数器移动到 \texttt{Circuit} 类后问题消失。
\end{attenzione}

\paragraph{步骤 4：下载。} 发送 HTTP 请求并一直读取到 \texttt{END}：

\begin{lstlisting}[language=Python]
s.send(b"GET /grande.bin HTTP/1.0\r\nHost: 127.0.0.1\r\n\r\n")
response = s.receive_all()
\end{lstlisting}

对一个 \num{1048576} 字节的随机文件，我们的一次运行得到：

\begin{lstlisting}[numbers=none]
流 1：209 字节，0.05 s；DATA 信元 2；截至目前 circuit SENDME 0；
          其他 ['END DONE']
流 2：1048781 字节，0.18 s；DATA 信元 2120；截至目前 circuit SENDME 68；
          其他 ['END DONE']
\end{lstlisting}

接收文件的 SHA-256 与原文件一致。68 条 \texttt{SENDME} 正好等于 $\lfloor(2+2120)/31\rfloor$：计数器确实属于整条电路。DATA 数量是 2120 而不是理论最小值 $\lceil1048781/498\rceil=2106$，因为出口不会一定等到凑满 498 字节才发送，而是按从服务器实际读取到的数据块进行封装。

\paragraph{步骤 5：观察窗口。} 接下来做一个普通客户端不会故意做的实验：下载同一个文件，但\emph{完全不发送} \texttt{SENDME}，观察出口发送多少 DATA 后停住；分别在启用和禁用拥塞控制时测试：

\begin{lstlisting}[language=Python]
circ = build(path, cc=cc)
circ.count_received_data = lambda: None      # 不发送 SENDME
\end{lstlisting}

\begin{lstlisting}[numbers=none]
congestion control=True：出口在 124 条 DATA 后停止
congestion control=False：出口在 500 条 DATA 后停止
\end{lstlisting}

启用拥塞控制时，124 正是初始拥塞窗口 $4\times31$。禁用时，出口在 500 而不是 1000 条停止，因为旧式流量控制下，\emph{流} window（500）先于 circuit window（1000）耗尽，客户端还必须发送流级 SENDME。两种情况下出口都只是等待，数据留在自身缓冲区中，而不是继续塞进 Tor 网络。

% ---------------------------------------------------------------------------
\section*{小结}
\begin{itemize}
  \item 流用 \texttt{BEGIN} 打开（\texttt{address:port} 必须小写），由 \texttt{CONNECTED} 确认，以 \texttt{DATA} 携带数据，以 \texttt{END} 关闭；客户端可以在 CONNECTED 前先发送乐观数据。
  \item 域名解析发生在出口，以避免向本地 DNS 解析器暴露目的地；\texttt{RESOLVE}/\texttt{RESOLVED} 可以在不打开连接的情况下解析名称。
  \item 原始流量控制使用每条电路 1000 信元、每条流 500 信元的窗口，并分别每 100/50 信元发送 SENDME。
  \item 电路级 \texttt{SENDME} 使用触发它的信元的累积摘要认证（\proposta{289}）。
  \item 使用拥塞控制时（\proposta{324}，默认 Vegas），窗口根据通过 SENDME 测得的 RTT 自适应；\file{sendme\_inc}=31，即一个 TLS 记录可装下的信元数，初始窗口 124。
  \item XON 与 XOFF 取代流窗口。
\end{itemize}

% ---------------------------------------------------------------------------
\section*{练习}

\begin{esercizio}[1]\label{es:11-begin}
为一个同时接受 IPv6 的客户端，写出连接 \texttt{Example.COM:443} 的 \texttt{BEGIN} 载荷的十六进制内容。
\end{esercizio}

\begin{esercizio}[1]\label{es:11-connected}
解释 \texttt{CONNECTED} 载荷 \texttt{c0000201 0000012c}。客户端应该怎样处理这些信息？
\end{esercizio}

\begin{esercizio}[1]\label{es:11-optimistic}
一条电路 RTT 为 600\,ms，一个页面只需要一次 HTTP 请求。估算 optimistic data 能节省多少时间。若连接最终失败，为什么出口必须丢弃已经排队的乐观数据？
\end{esercizio}

\begin{esercizio}[2]\label{es:11-velocita}
在原始流量控制下，一条 RTT 为 400\,ms 的电路上，单条流的最大速度大约是多少？若同一电路有多条流呢？把这些数值与拥塞控制的初始窗口比较。
\end{esercizio}

\begin{esercizio}[2]\label{es:11-sendme-falso}
描述一种利用未认证 \texttt{SENDME} 耗尽出口内存的攻击。为什么 DATA 信元中的随机填充对认证有效性很重要？
\end{esercizio}

\begin{esercizio}[2]\label{es:11-resolve}
为 \file{stream.py} 增加 \texttt{resolve(name)}：发送 \texttt{RESOLVE} 并解析 \texttt{RESOLVED}。在实验网络中用 \texttt{localhost} 测试。不存在的名称会让出口怎样回复？
\end{esercizio}

\begin{esercizio}[3]\label{es:11-policy}
修改实验网络出口 \texttt{test006e} 的 \file{torrc} 出口策略并执行 \cmd{chutney hup}，然后尝试打开一个被拒绝端口。\texttt{END} 返回哪个原因码？解释后续附加字段。
\end{esercizio}

\begin{esercizio}[3]\label{es:11-rtt}
修改实验代码，在客户端测量发送每条 \texttt{SENDME} 到收到下一条 \texttt{DATA} 的时间。你实际测量的是什么？再在同一条电路并行进行第二个下载，比较结果。
\end{esercizio}

% ---------------------------------------------------------------------------
\section*{延伸阅读}
\begin{itemize}
  \item \titolo{Tor Protocol Specification} 中的 \enquote{Opening streams and transmitting data}、\enquote{Closing streams}、\enquote{Remote hostname lookup} 与 \enquote{Flow control}~\cite{torspec}。
  \item 提案 289，\titolo{Authenticating sendme cells to mitigate bandwidth attacks}。
  \item 提案 324，\titolo{RTT-based Congestion Control for Tor}：设计动机、算法、参数与安全分析。
\end{itemize}
